\documentclass[a4paper,11pt]{article}
\usepackage{exam-style-842-unit}
\usepackage{needspace}
\examinehead{842 信息化概论 · 第06单元练习册}{练习二\quad 传感器、RFID 与传感器网络}
\begin{document}

\unitcover{842 信息化概论 · 单元练习册}{第06单元\quad 物联网与传感技术}{练习二\quad 传感器、RFID 与传感器网络（共 100 分）}{约 75 分钟}{微电子技术概述；传感器的定义、组成（敏感元件·转换元件·转换电路·辅助电源）、作用与分类；静态特性 8 指标与动态特性；数据采集原理；智能传感器与 MEMS；RFID 系统组成、工作原理及与条码的比较；无线传感器网络的节点、自组网、路由与能耗。}

\begin{knowmap}
\textbf{微电子基础：}微电子技术以集成电路为核心，历经电子管、晶体管、集成电路、集成电路规模化、微机电集成电路五代；是传感、通信、计算机技术发展的重要基础\\[2pt]
\textbf{传感器与采集链：}能感受规定的被测量并按一定规律转换成可用信号的器件或装置（GB7665-87），由敏感元件＋转换元件组成（实际含转换电路、辅助电源）；被测量 $\rightarrow$ 敏感元件（预变换）$\rightarrow$ 转换元件 $\rightarrow$ 转换电路 $\rightarrow$ 电量输出\\[2pt]
\textbf{分类与特性：}按被测物理量·工作原理·输出信号性质·能量传递方式分类；静态特性 8 指标＝量程·灵敏度·线性度·不重复性·迟滞·精确度·零漂·温漂；动态特性＝频率响应＋瞬态响应，常简化为一阶或二阶系统\\[2pt]
\textbf{RFID：}射频识别，非接触式自动识别技术，系统由电子标签、读写器、天线和后端信息系统组成；标签进入读写器射频场获能并与读写器交换信息\\[2pt]
\textbf{无线传感器网络：}节点＝传感器模块＋处理器模块＋无线通信模块＋能量供应模块；自组织、多跳、以数据为中心、能量受限\\[2pt]
\textbf{MEMS：}微机电系统，集微型传感器、微型执行器、信号处理器与控制电路、接口电路、通信系统及电源于一体，特点为微型化、集成化、低成本
\end{knowmap}

\begin{keyterms}
\textbf{传感器}：能感受规定的被测量并按一定规律转换成可用信号的器件或装置，是把外界非电量信息转换成电量输出的器件。\\
\textbf{敏感元件}：能直接感受被测量（一般为非电量）并输出与被测量成确定关系的其他物理量的元件，又称预变换器。\\
\textbf{转换元件}：直接或不直接感受被测量，并把敏感元件的输出量转换成电量后再输出的元件，又称传感元件。\\
\textbf{转换电路}：把转换元件输出的电参量转换成电压、电流或频率量的电路，如电阻应变片传感器采用的电桥测量电路。\\
\textbf{辅助电源}：为需要电源才能工作的转换电路和转换元件提供正常工作电源，通常采用直流电源。\\
\textbf{灵敏度}：稳态时输出量变化量与输入量变化量的比值，即单位输入量引起的输出量大小，比值越大传感器越灵敏。\\
\textbf{线性度}：传感器校准曲线与选定的拟合直线的偏离程度，又称非线性误差。\\
\textbf{迟滞}：传感器正向特性曲线（输入量增大）与反向特性曲线（输入量减小）的不一致程度。\\
\textbf{精确度}：又称精度，是线性度、不重复性及迟滞 3 项指标的综合指数，反映系统误差和随机误差的综合指标。\\
\textbf{智能传感器}：由传统传感器与微处理器（或微计算机）相结合构成，具有检测、判断、信息处理、信息记忆、逻辑思维等功能。\\
\textbf{RFID}：射频识别技术，通过射频信号非接触地自动识别目标对象并获取相关数据，是物联网中让物品“开口说话”的关键技术。\\
\textbf{电子标签}：RFID 系统中附着在物品上、存储规范而具有互用性信息的部分，分无源标签和有源标签。\\
\textbf{无线传感器网络}：由大量传感器节点以自组织、多跳方式构成的无线网络，协作感知、采集和处理覆盖区域内被监测对象的信息。\\
\textbf{MEMS}：微机电系统，集微型传感器、微型执行器、信号处理器和控制电路、接口电路、通信系统以及电源于一体的微型机电系统。
\end{keyterms}

\needspace{6cm}
\begin{contrast}{传感器的组成（每行依次为：定义·作用·典型举例）}
\citem{敏感元件}{能直接感受被测量并输出与被测量成确定关系的其他物理量的元件}{完成预变换，把难以直接转换的非电量变成易于转换的非电量}{压力传感器中的膜片先把压力转换为位移}
\citem{转换元件}{把敏感元件的输出量转换成电量后再输出}{完成非电量到电量的变换}{把位移量变换为电容量、电阻量或电感量}
\citem{转换电路}{把转换元件输出的电参量转换成电压、电流或频率量的电路}{把电参量变成便于传输处理的电信号}{电阻应变片传感器采用的电桥测量电路}
\citem{辅助电源}{为需要电源才能工作的转换电路和转换元件供电}{保证电路正常工作}{通常采用直流电源}
\end{contrast}

\needspace{6cm}
\begin{contrast}{传统传感器 / 智能传感器（每行依次为：组成·核心部件·功能特点）}
\citem{传统传感器}{敏感元件与转换元件（含转换电路、辅助电源）}{敏感元件与转换元件}{只完成被测量到电量的转换，功能单一，需外部标定}
\citem{智能传感器}{传统传感器与微处理器（或微计算机）相结合}{微处理器}{具有自补偿、自检自校自诊断、复合敏感、通信接口、现场学习、掉电保护等功能}
\end{contrast}

\needspace{6cm}
\begin{contrast}{RFID / 条形码（每行依次为：识别方式·信息可改写性·读取特点）}
\citem{RFID}{非接触式，通过射频信号自动识别}{信息可重复读写、可加密}{可批量读取、作用距离从几厘米到几十米、抗污损能力强}
\citem{条形码}{接触式或近距离光电扫描识读}{信息一般印刷后不可改写}{需逐件近距离扫描、易磨损污染、成本低}
\end{contrast}

\needspace{6cm}
\begin{contrast}{无线传感器网络 / 传统有线监测系统（每行依次为：组网方式·部署与扩展·关键约束）}
\citem{无线传感器网络}{节点自组织成多跳无线网络，无需固定基础设施}{部署灵活、可大规模密集布设、扩展方便}{节点电池供电，能量受限是设计核心约束}
\citem{传统有线监测系统}{传感器通过线缆点对点接入采集设备}{布线成本高、部署位置受限、扩展困难}{供电与布线受现场条件限制}
\end{contrast}

\qtypename{一、单项选择题}{10 题，每题 3 分，共 30 分}
\q 微电子技术的核心是（\quad）
\choicesline{A. 集成电路技术}{B. 传感器技术}{C. 通信技术}{D. 计算机技术}

\q 能感受规定的被测量并按照一定规律转换成可用信号的器件或装置是（\quad）
\choicesline{A. 传感器}{B. 执行器}{C. 控制器}{D. 显示器}

\q 压力传感器中的膜片先把压力转换为位移，该膜片属于（\quad）
\choicesline{A. 敏感元件}{B. 转换元件}{C. 转换电路}{D. 辅助电源}

\q 传感器的静态特性指标中，单位输入量引起的输出量大小称为（\quad）
\choicesline{A. 量程}{B. 灵敏度}{C. 线性度}{D. 迟滞}

\q 传感器正向特性曲线（输入量增大）与反向特性曲线（输入量减小）的不一致程度称为（\quad）
\choicesline{A. 不重复性}{B. 线性度}{C. 迟滞}{D. 零漂}

\q 线性度、不重复性及迟滞 3 项指标的综合称为（\quad）
\choicesline{A. 灵敏度}{B. 精确度}{C. 分辨率}{D. 量程}

\q 智能传感器的核心是（\quad）
\choicesline{A. 敏感元件}{B. 微处理器}{C. 转换电路}{D. 辅助电源}

\q RFID 系统中附着在物品上、存储规范而具有互用性信息的部分是（\quad）
\choicesline{A. 电子标签}{B. 读写器}{C. 天线}{D. 后端信息系统}

\q 无线传感器网络节点通常由电池供电，其设计中最突出的约束是（\quad）
\choicesline{A. 能量}{B. 存储容量}{C. 显示分辨率}{D. 计算速度}

\q 1947 年合作发明晶体管的贝尔实验室三位科学家是（\quad）
\choices{A. 肖克利、巴丁、布拉顿}{B. 基尔比、诺伊斯、摩尔}
{C. 香农、维纳、贝塔朗菲}{D. 盖茨、乔布斯、艾伦}

\qtypename{二、多项选择题}{4 题，每题 4 分，共 16 分；每题有两个或两个以上正确项}
\q 关于传感器的组成，下列说法正确的有（\quad）
\choices{A. 一般由敏感元件和转换元件组成}{B. 实际应用还包括转换电路}
{C. 辅助电源为需要电源的部件供电}{D. 转换电路是任何传感器都必须配置的}

\q 传感器常见的分类方法有（\quad）
\choicesline{A. 按被测物理量}{B. 按工作原理}{C. 按输出信号性质}{D. 按生产厂家规模}

\q 与条形码相比，RFID 的优点有（\quad）
\choices{A. 非接触识别}{B. 信息可重复读写}{C. 可批量读取多个标签}{D. 标签成本更低}

\q 无线传感器网络的主要特点有（\quad）
\choicesline{A. 自组织组网}{B. 多跳路由传输}{C. 以数据为中心}{D. 节点能量不受限制}

\qtypename{三、判断题}{6 题，每题 2 分，共 12 分；正确的画“√”，错误的画“×”并写出改正后的表述}
\q 微电子技术以集成电路为核心，至今已经历了五代技术的突破与发展。\quad（\quad）\par\noindent 改正：\fixblank

\q 传感器一般由敏感元件和转换元件两个基本元件组成，实际应用中还包括转换电路和辅助电源。\quad（\quad）\par\noindent 改正：\fixblank

\q 传感器在相同条件下按同一方向做全量程多次重复测量时输出曲线的不一致程度称为迟滞。\quad（\quad）\par\noindent 改正：\fixblank

\q 智能传感器是由传统传感器与微处理器（微计算机）相结合构成的，微处理器是其核心。\quad（\quad）\par\noindent 改正：\fixblank

\q RFID 属于接触式自动识别技术，必须贴近读写器才能读取标签信息。\quad（\quad）\par\noindent 改正：\fixblank

\q 无线传感器网络节点通常由电池供电、能量不受限制，因此可以长期持续高速传输数据。\quad（\quad）\par\noindent 改正：\fixblank

\qtypename{四、填空题}{6 空，每空 1 分，共 6 分}
\q 传感器中能直接感受被测量并输出与被测量成确定关系的其他物理量的元件称为\fillblanks[2.0cm]。

\q 灵敏度是稳态时输出量变化量与\fillblanks[2.6cm]的比值。

\q 传感器静态特性中，线性度、不重复性及迟滞 3 项指标的综合称为\fillblanks[1.8cm]。

\q RFID 系统一般由电子标签、\fillblanks[1.8cm]、天线和后端信息系统组成。

\q 无线传感器网络的英文缩写是\fillblanks[1.2cm]，其节点通常由电池供电，能量受限是网络设计的核心约束。

\q MEMS 的中文名称是\fillblanks[2.8cm]。

\newpage
\qtypename{五、简答题}{4 题，每题 6 分，共 24 分}
\q 什么是传感器？简述传感器的一般组成及各部分的作用。\anslines{2.4cm}

\q 简述传感器静态特性的 8 个主要指标及其含义。\anslines{2.4cm}

\q 什么是 RFID？简述 RFID 系统的组成与基本工作原理，并说明其与条形码相比的主要优点。\anslines{2.4cm}

\q 什么是无线传感器网络？简述其节点组成、主要特点以及能量受限条件下的节能思路。\anslines{2.4cm}

\newpage
\qtypename{六、综合题}{1 题，共 12 分}
\q 某奶牛场为每头奶牛佩戴电子耳标，在牛舍安装温湿度传感器和氨气传感器，在挤奶厅入口设置 RFID 读写器，数据经 ZigBee、4G 传至云端管理平台。请回答：（1）以氨气传感器为例，说明传感器数据采集的完整链条（被测量到电量输出）；（2）指出该场 RFID 系统由哪些部分组成，并说明无源电子耳标如何获得工作能量；（3）分析电子耳标相比传统条形码耳标的优势，说明 MEMS 传感器的主要特点，并给出 3 条降低无线传感器网络能耗的具体措施。\anslines{5cm}

\end{document}
